% %--------------------------
% %	CONFIGURATION
% %--------------------------
\documentclass[11pt,a4paper]{moderncv}
\moderncvstyle{classic}
\moderncvcolor{black}
\definecolor{firstnamecolor}{rgb}{0,0,0}

% Basic packages
\usepackage[utf8]{inputenc}
\usepackage[T1]{fontenc}
\usepackage[scale=0.8, top=1.5cm, bottom=1.5cm]{geometry}
\usepackage[dvipsnames]{xcolor}
\definecolor{PaleBlue}{RGB}{76, 151, 215}
\usepackage[colorlinks=true, urlcolor=PaleBlue, linkcolor=PaleBlue]{hyperref}
\usepackage{microtype}

% Configure font settings
\renewcommand{\rmdefault}{ppl}
\renewcommand{\sfdefault}{phv}
\renewcommand*{\namefont}{\fontsize{26}{18}\mdseries}

% Configure microtype
\microtypesetup{expansion=false}

% Load sensitive information
\input{../../secrets/personal_info_dk.tex}

% %--------------------------
\begin{document}
\makecvtitle
Dear Members of the Selection Committee,

\vspace*{2em}
I am applying for the PhD Fellowship in Simulation-Supervised Machine Learning for Biology at the University of Copenhagen. After seven years working with machine learning in research and industry, I want to return to research. I enjoy working on difficult scientific problems with people I can learn from. I wish to pursue problems that require time and dedication to investigate why a method works, where it fails, and how it can be improved. 

\vspace{4pt}
\hspace*{2em}
This project attracts me as it offers the opportunity to return to research in machine learning for medical imagery. I would like to investigate how simulations can be tuned to capture the variation and noise in real images, and how to assess whether models trained on them generalise to experimental data. The combination of differentiable simulation and image analysis would allow me to build on my background in both physics and machine learning.

\vspace{4pt}
\hspace*{2em}
At RaySearch Laboratories, I researched methods for automated radiotherapy planning, first through my master's thesis and then as a full-time researcher. I developed anatomical representations from segmented CT scans using autoencoders and used sparse Gaussian processes to predict dose-volume histograms and spatial dose distributions. I also worked on lexicographic optimisation methods to better represent clinical priorities. I enjoyed the technical work and environment equally: discussing difficult questions with medical physicists and researchers, testing ideas, and learning enough about the underlying problem to make useful contributions. In addition, improving cancer treatment gave the work a purpose.

\vspace{4pt}
\hspace*{2em}
Since then, I have gained experience taking responsibility for larger computational systems. At Valcon, I led machine learning projects from problem formulation through to implementation and deployment. At Signum Life Science, I now work on the architecture and development of a forecasting platform supporting thousands of independently trained models. These roles have taught me to structure complex software, design evaluation procedures, investigate failures, and work independently while coordinating with others. Independently, I also built Iris, a computer-vision retrieval system combining object localisation, image embeddings, and similarity search. I have included it as a code sample with a short project description.


\vspace{4pt}
\hspace*{2em}
My strongest experience is in machine learning, mathematical modelling, and Python development. I would need to deepen my knowledge of biophysical simulation and microscopy, but my training in physics and mathematics, combined with my experience learning new domains, gives me confidence that I can do so. My degree in Engineering Physics, with a specialisation in Statistical Learning and Data Analytics, gives me a strong foundation to build on. I would bring experience developing and evaluating methods, building the software needed to use them, and working with specialists whose knowledge complements my own.

\vspace*{2em}
I welcome the opportunity to discuss the project and where my background could be most useful.

\vspace{2em}
Kind regards, 

Ivar Cashin Eriksson

\end{document}
